\section{调和分析：变换的两端与归一化的位置}

\begin{know}{H1}{DFT 的定义与「正向不归一化」}{框架：math/fourier.hpp}
\field{是什么} $X_k=\sum_{n=0}^{N-1}x_n e^{-2\pi ikn/N}$，$k=0,\dots,N-1$；正向\textbf{不除 $N$}。

\field{为什么需要} 它是「把一个离散信号分解成 $N$ 个旋转的叠加」；不归一化让正向成为纯投影，
所有归一化集中到逆变换，避免因子分散在两处。

\field{关键结论} 旋转因子 $e^{-2\pi ikn/N}$ 对应 C3 的复指数；$X_k$ 是 $x$ 在频率 $k$ 上的\textbf{内积}
$\inner{e_k}{x}$（严格说要共轭，方向由符号约定决定）。成立条件：$N\ge1$；信号按 $N\times1$ 列向量（D19）。

\field{判据} 与朴素 DFT 对拍：$N=1024$ 复数随机，$\norm{\mathrm{FFT}-\mathrm{DFT}}_{\infty}/\norm{X}_{\infty}\le10^{-13}$；
另加\textbf{有条件逐点}判据——只对 $\abs{X_k}\ge0.1\norm{X}_{\infty}$ 的频点比较（近零频点必然假红）。

\field{前置} C3、L4。

\field{常见误解} 以为「不归一化」是无所谓的选择。它决定了卷积定理里有没有附加因子（H3）、
Parseval 里 $1/N$ 放在哪边（H8）——一处改，处处变。
\end{know}

\begin{know}{H2}{逆变换：$1/N$ 只出现一次}{框架：math/fourier.hpp}
\field{是什么} $x_n=\frac1N\sum_k X_k e^{+2\pi ikn/N}$。

\field{为什么需要} 正逆必须复合为恒等：$\frac1N\sum_k e^{2\pi ik(n-m)/N}=\delta_{nm}$。
这个 $1/N$ 是「$N$ 个正交方向的平均」的代价。

\field{关键结论} 两个方向的\textbf{符号相反}是定义本身，不是笔误；$1/N$ 全库只在 IFFT 出现一次。
成立条件：$N$ 与正向一致；复数域共轭关系 $x=\overline{\mathrm{FFT}(\overline{X})}/N$ 可作为实现路径。

\field{判据} 往返 $\norm{\mathrm{IFFT}(\mathrm{FFT}(x))-x}_{\infty}\le10^{-13}$。

\field{前置} H1。

\field{常见误解} 以为「正变换除 $\sqrt N$、逆变换除 $\sqrt N$」更对称。那是酉归一化（Parseval 无因子），
与 D20 约定的 $\pm1$、$1/N$ 方案不同；混用会让卷积少（或多）一个因子。
\end{know}

\begin{know}{H3}{循环卷积定理}{框架：math/fourier.hpp}
\field{是什么} $(a\circledast b)_n=\sum_{m}a_m b_{(n-m)\bmod N}$，则
$\mathrm{FFT}(a\circledast b)_k=A_kB_k$（\textbf{无附加因子}）。

\field{为什么需要} 时域卷积变成频域乘法，是 FFT 全部实用价值的来源。

\field{关键结论} 前提是\textbf{循环}卷积（下标模 $N$）；无附加因子正是不归一化 + $1/N$ 只在逆变换的后果。
成立条件：两序列等长 $N$；不等长要先补零到同一长度，否则得到的是循环的、绕回的结果。

\field{判据} 三条路径对拍：时域直接、频域补零、循环矩阵乘，逐点 $\le10^{-12}$。

\field{前置} H1、H2。

\field{常见误解} 以为「FFT 卷一下就得到线性卷积」。循环卷积会把尾部\textbf{绕回}头部；
结果看起来「差不多」但边界全错，这是最常见的静默错。
\end{know}

\begin{know}{H4}{线性卷积与补零长度}{框架：math/fourier.hpp}
\field{是什么} $(a*b)_n=\sum_m a_mb_{n-m}$，长度 $N+M-1$；用 FFT 时补零到
$L\ge N+M-1$ 且 $L=2^{k}$。

\field{为什么需要} 线性卷积是滤波、多项式乘法的标准形式；只有补够零，循环卷积才与它相等。

\field{关键结论} $L<N+M-1$ 时尾部环绕（$y_{n}$ 与 $y_{n+L}$ 混在一起），这是\textbf{另一个函数}而非精度问题。
源码取 $L=\text{NextPowerOfTwo}(N+M-1)$，且\textbf{不做自动分派}——非 2 的幂不会偷偷退回 $O(N^{2})$ 的 DFT。

\field{判据} 前 $N+M-1$ 项与直接卷积逐点 $\le10^{-12}$；$L$ 取小一位时应当出现可见差异（反例用例）。

\field{前置} H3。

\field{常见误解} 以为「补零只是形式」。补零长度决定\textbf{语义}：够长才是线性卷积，
不够长就是循环卷积，两者在边界差一整段。
\end{know}

\begin{know}{H5}{频轴与 Nyquist}{框架：math/fourier.hpp}
\field{是什么} 采样间隔 $dx$ 下，$\xi_k=k/(N\,dx)$（$k\le N/2$）；后半段
$\xi_k=(k-N)/(N\,dx)$（$k>N/2$）。采样率 $f_s=1/dx$，Nyquist 频率 $f_s/2$ 落在 $k=N/2$。

\field{为什么需要} 没有频轴，$X_k$ 只是「第 $k$ 个数」；有了频轴才能说「这是 3 Hz 的分量」。

\field{关键结论} 前半段是正频率、后半段是负频率（因为 DFT 是周期谱，$\xi$ 与 $\xi+f_s$ 不可分辨）。
成立条件：等距采样；$dx>0$。$N$ 偶数时 $k=N/2$ 是自共轭点（实信号下该点相位移 $\pi$ 无意义）。

\field{判据} 直流分量在 $k=0$；实信号在 $k$ 与 $N-k$ 共轭（H8 的对称性）。

\field{前置} H1。

\field{常见误解} 以为频点顺序是「从 0 到最大频率」。DFT 的排列是
$0,1,\dots,N/2,-N/2+1,\dots,-1$，中间那个折返点最容易被忽略。
\end{know}

\begin{know}{H6}{radix-2：位反转与蝶形}{框架：math/fourier.hpp}
\field{是什么} 按频率下标的二进制位反转做一次重排，然后做 $\log_2N$ 层蝶形，每层步长 $2,4,\dots,N$。

\field{为什么需要} 把 $O(N^{2})$ 的求和拆成两半递归，是「分治」在变换上的原型；
它也是「算法结构决定常数」的样板。

\field{关键结论} 位反转的原因：递归拆偶/奇下标后，叶子顺序是「下标反转二进制」。
$N=1$ 合法（恒等）。成立条件：$N=2^{k}$；旋转因子必须\textbf{直接算 $\cos/\sin$ 建表}——
用递推 $w^{k+1}=w^{k}\cdot w$ 会把误差累到 $N\eps$（$N=2^{20}$ 时约 $2\times10^{-10}$，跨过 $10^{-12}$ 判据）。

\field{判据} 与 DFT 对拍（H1）；$\log$-$\log$ 斜率不是判据（这是精确算法，不是逼近）。

\field{前置} H1、A2。

\field{常见误解} 以为「FFT 是一种近似」。除了浮点舍入，FFT 与 DFT 给出\textbf{完全相同}的结果；
速度来自因子分解，不来自截断。
\end{know}

\begin{know}{H7}{复杂度：$O(N^{2})$ 与 $O(N\log N)$ 的量级差}{框架：math/fourier.hpp}
\field{是什么} DFT 需 $N^{2}$ 次复乘；FFT 需约 $\frac N2\log_2N$ 次蝶形（每次一个复乘两个复加）。

\field{为什么需要} 它解释「为什么 FFT 是信号处理的起点」：$N=2^{20}$ 时约 $10^{12}$ 对 $2\times10^{7}$，
差四个数量级。

\field{关键结论} 比值约 $2N/\log_2N$。成立条件：$N$ 为 2 的幂；常数因子未计。

\field{判据} 无（复杂度是量级陈述，不进 ctest）；实测只记一笔不设阈值。

\field{前置} H6。

\field{常见误解} 以为「$O(N\log N)$ 一定比 $O(N^{2})$ 快」。小 $N$（几十）时 DFT 的常数更小，
这也是库在小规模时用 DFT 作参照基线的原因。
\end{know}

\begin{know}{H8}{Parseval 与归一化因子}{框架：math/fourier.hpp}
\field{是什么} 在不归一化 + $1/N$ 逆变换的约定下：$\sum_n\abs{x_n}^{2}=\frac1N\sum_k\abs{X_k}^{2}$。

\field{为什么需要} 能量守恒是检验「变换有没有算错」的全局判据——它不依赖任何单个频点的精度。

\field{关键结论} 因子 $1/N$ 来自 H2 的逆变换归一化；若改用酉归一化（正逆各除 $\sqrt N$），因子变成 1。
成立条件：同一约定下两边比较；$\mathrm{abs2}$ 求和（不取模开方）。

\field{判据} 两边比值与 1 的差 $\le10^{-12}$。

\field{前置} H1、H2。

\field{常见误解} 以为 Parseval 是「近似成立的能量关系」。它是精确恒等式（离散、有限、无截断），
浮点上只差舍入。
\end{know}

\begin{know}{H9}{傅里叶级数：系数就是归一化谱}{框架：math/series.hpp}
\field{是什么} $c_k=\frac1N\sum_{n=0}^{N-1}x_ne^{-2\pi ikn/N}=X_k/N$，$k\in\Z$；
布局为 $(2K+1)\times1$，第 $r$ 行对应 $k=r-K$。

\field{为什么需要} 把「周期函数的频谱」与「离散 DFT」接上：级数系数不是新对象，就是 DFT 除 $N$。

\field{关键结论} 手算锚点：$x_n=\cos(2\pi\cdot3n/N)\Rightarrow c_{\pm3}=1/2$，其余 $0$。
截断重构 $f_K(t)=\sum_{k=-K}^{K}c_ke^{2\pi ikt}$ 在节点上等于 $N\cdot\mathrm{IFFT}(S)$——
\textbf{乘 $N$ 不是笔误}，因为 IFFT 自带 $1/N$（H2）。成立条件：$N=2^{k}$、$0\le K$、$2K+1\le N$；
截断窗口必须\textbf{对称}（单边截断得到的是解析信号，实部只剩一半）。

\field{判据} 漏乘 $N$ 的病态要能被幅度判据抓住（$f_K$ 与实际值差一个 $N$ 倍）。

\field{前置} H1、H2。

\field{常见误解} 以为 $c_k$ 与 $X_k$ 是「两个不同的量」。差别只有一个常数 $1/N$；
真正的差别在\textbf{索引约定}（$k$ 取负、居中存放）。
\end{know}

\begin{know}{H10}{截断与 Gibbs 过冲}{框架：math/series.hpp}
\field{是什么} 用前 $K$ 个谐波重构有跳变的周期函数时，跳变附近出现过冲，峰值趋于
$\frac2\pi\int_0^{\pi}\frac{\sin t}{t}dt=1.08948987\ldots$，即约 $8.95\%$ 的相对过冲。

\field{为什么需要} 它是「逐点收敛」与「一致收敛」分家的经典例子；也是「误差不随 $K$ 减小」的唯一常见例外。

\field{关键结论} 过冲幅度\textbf{与 $K$ 无关}（只与跳变高度有关），过冲宽度 $\sim1/K$——
所以 $L^{2}$ 误差趋于 0（H11），而 $\sup$ 误差不趋于 0。成立条件：函数有第一类间断。

\field{判据} 方波重构的峰值比 $\approx1.0895$，且 $K$ 增大时该比值不降；「不随 $K$ 减小」本身就是判据。

\field{前置} H9。

\field{常见误解} 以为「加更多项就能消除振铃」。可以把振铃\textbf{挤窄}，不能压平；
要压平必须换求和法（Fejér / Lanczos $\sigma$ 因子），那是改变函数而非提高分辨率。
\end{know}

\begin{know}{H11}{谱收敛：光滑度决定衰减速度}{框架：math/series.hpp}
\field{是什么} 若周期函数 $x$ 有 $p$ 阶连续导数（$x\in C^{p}$），则 $\abs{c_k}=O(\abs{k}^{-p-1})$。

\field{为什么需要} 它把「函数有多光滑」翻译成「系数掉得多快」，是选 $K$ 的唯一依据。

\field{关键结论} 方波 $p=0\Rightarrow\abs{c_k}=O(1/\abs{k})$，截断误差 $O(1/K)$；
解析函数则指数衰减。成立条件：$p$ 阶导数连续（分段光滑时按最差的那段算）。

\field{判据} 双对数坐标下 $\abs{c_k}$ 的斜率约 $-(p+1)$。

\field{前置} H9。

\field{常见误解} 以为「光滑函数的系数一定很快衰减」。$C^{1}$ 但二阶导数跳变的函数，
衰减只有 $O(k^{-3})$；指数衰减要求\textbf{解析}，比「无限可微」更强。
\end{know}
